\documentclass[11pt]{article}
\usepackage[a4paper,margin=1.45cm]{geometry}
\usepackage[T1]{fontenc}
\usepackage[utf8]{inputenc}
\usepackage{lmodern}
\usepackage[spanish,es-noquoting,es-noshorthands]{babel}
\usepackage{amsmath,amssymb,mathtools,array,booktabs}
\usepackage[protrusion=true,expansion=false]{microtype}
\setlength{\parindent}{1.25em}
\setlength{\parskip}{0pt}
\allowdisplaybreaks
\setlength{\emergencystretch}{30em}
\sloppy
\hbadness=10000
\vbadness=10000
\hfuzz=10pt
\vfuzz=10pt
\raggedbottom
\providecommand{\mG}{\mathfrak{G}}
\providecommand{\mA}{\mathfrak{A}}
\providecommand{\mB}{\mathfrak{B}}
\providecommand{\mC}{\mathfrak{C}}
\providecommand{\mD}{\mathfrak{D}}
\providecommand{\mE}{\mathfrak{E}}
\providecommand{\mH}{\mathfrak{H}}
\providecommand{\mK}{\mathfrak{K}}
\providecommand{\mM}{\mathfrak{M}}
\providecommand{\mN}{\mathfrak{N}}
\providecommand{\mR}{\mathfrak{R}}
\providecommand{\mS}{\mathfrak{S}}
\providecommand{\mT}{\mathfrak{T}}
\providecommand{\mU}{\mathfrak{U}}
\providecommand{\mO}{\mathfrak{O}}
\providecommand{\mo}{\mathfrak{o}}
\providecommand{\OO}{\Omega}
\providecommand{\Th}{\Theta}
\providecommand{\iso}{\simeq}
\providecommand{\normlhd}{\triangleleft}
\newcommand{\mat}[1]{\begin{pmatrix}#1\end{pmatrix}}
\begin{document}

\setcounter{footnote}{0}

\section*{34. Magnitudes hipercomplejas y teoría de representaciones}
\emph{Math. Zs. 30 (1929), pp. 641--692.}

\subsection*{Introducción}

Los teoremas generales más importantes sobre los sistemas hipercomplejos se remontan a Molien (\emph{Math. Annalen}, vol. 41, y los informes de Dorpat de 1897). De modo esencialmente independiente, Frobenius desarrolló poco después, como una teoría unificada, la teoría de los sistemas hipercomplejos y de sus representaciones, en particular la teoría de representaciones de los grupos finitos. El fundamento es el concepto, debido a Dedekind, del determinante de grupo y, más en general, del determinante de un sistema hipercomplejo arbitrario. Frobenius muestra que a los distintos factores irreducibles de este determinante (tomando siempre como dominio de coeficientes el cuerpo de todos los números complejos) corresponden las distintas clases de representaciones irreducibles, y que de este modo se agotan todas las clases de representaciones irreducibles. Tal clase de representación está completamente caracterizada por su ``sistema de caracteres''; en el caso de los grupos finitos, este sistema de caracteres se obtiene al descomponer el determinante del sistema hipercomplejo conmutativo que se deriva de las clases de elementos conjugados del grupo. Este determinante se descompone en factores lineales, con los caracteres como coeficientes: una generalización directa del resultado obtenido primero por Dedekind, a saber, que el determinante de grupo de un grupo abeliano finito se descompone en factores lineales cuyos coeficientes son los distintos caracteres del grupo abeliano (correspondencia con Frobenius).\footnote{Lo que sigue es una elaboración libre, preparada por B. L. van der Waerden, de mis lecciones del semestre de invierno de 1927/28. Preparamos conjuntamente el texto para la imprenta. También debo a B. L. van der Waerden varias observaciones críticas.}

Estos resultados, conceptualmente sencillos y transparentes, son obtenidos sin embargo por Frobenius mediante cálculos laboriosos. El desarrollo posterior buscó una derivación simplificada de esos resultados y, al mismo tiempo, su extensión cuando se toma como dominio de coeficientes un cuerpo arbitrario. Ese desarrollo posterior avanzó por caminos completamente separados para los sistemas hipercomplejos y para la teoría de representaciones. Los sistemas hipercomplejos recibieron un tratamiento aritmético por Wedderburn: todo sistema hipercomplejo es, de manera única, una suma directa de anillos bilateralmente directamente indescomponibles; para los sistemas sin radical, esos anillos indescomponibles resultan isomorfos al sistema de todas las matrices de $n$ filas con elementos de un anillo de división asignado, no necesariamente conmutativo, determinado esencialmente de manera única.\footnote{Cf. las referencias en Dickson, \emph{Algebras and their Arithmetics} (edición alemana: \emph{Algebren und ihre Zahlentheorie}, Zúrich 1927).} La teoría de representaciones recibió una fundamentación elemental por Burnside e I. Schur, quienes partieron directamente de una representación dada, independientemente del sistema hipercomplejo, y operaron con teoremas de matrices. En particular Schur trató la cuestión de los cuerpos numéricos de menor grado en los que una representación irreducible sobre un cuerpo dado se descompone absolutamente. Estos componentes absolutamente irreducibles pertenecen a un número finito de clases de representaciones conjugadas. El grado de los cuerpos numéricos buscados es igual al producto del número de esas clases por el índice (el número de componentes en cada una de las clases conjugadas). El principal instrumento auxiliar de la teoría de Schur es el sistema de todas las matrices que conmutan con una representación irreducible.\footnote{Cf. las referencias en el capítulo 10 de la teoría de grupos de Speiser.}

En la fundamentación puramente aritmética que sigue, las magnitudes hipercomplejas y la teoría de representaciones aparecen de nuevo como un todo unificado: como caso particular de una teoría general de anillos no conmutativos que satisfacen solo ciertas condiciones de finitud. Más precisamente, se trata de una teoría de clases de módulos y de ideales respecto de esos anillos, con el resultado principal de que las clases de módulos irreducibles ya quedan agotadas por las clases de ideales correspondientes; en particular, para los anillos sin radical, todas las clases de módulos se descomponen en irreducibles (son completamente reducibles). Este es el equivalente aritmético del resultado de Frobenius según el cual los factores irreducibles del determinante regular de grupo, respectivamente de sistema, agotan ya la totalidad de las clases de representaciones irreducibles, y según el cual la reducibilidad completa de las representaciones vale para los sistemas sin radical. En efecto, la consideración del determinante del sistema y de su descomposición puede interpretarse como un paso a la norma, mientras que aquí se tratan directamente la descomposición en suma directa de los ideales y sus series de composición. La teoría de representaciones, mediante el concepto de módulo de representación, se convierte en una teoría de clases de módulos.

La introducción de las clases de módulos y de ideales tiene una base puramente teórica de grupos. Los módulos y los ideales se consideran como grupos abelianos respecto de la adición, restringidos por la condición de que admitan ciertas multiplicaciones por elementos del anillo: forman ``grupos con operadores'' (§1). Para tales grupos, en lugar de la isomorfía ordinaria aparece la ``isomorfía por operadores'' (§2); los grupos que son operador-isomorfos entre sí (dentro de un dominio fijo) se reúnen en una clase, concepto de clase que, por ejemplo en el caso de los ideales de un cuerpo numérico, coincide con el usual. La teoría de los grupos con operadores se remonta a W. Krull y O. Schmidt (cf. nota 6); se desarrolla sistemáticamente en el capítulo I. Los teoremas que siguen siendo válidos aquí también -- sobre series de composición, producto directo (respectivamente suma), y grupos completamente reducibles; teoremas de unicidad en el sentido de la división en clases recién mencionada -- forman la base de todo lo que sigue.

De allí resultan primero los teoremas generales de unicidad para los constituyentes diagonales irreducibles de representaciones arbitrarias (capítulo III, §16), sobre la base de la correspondencia biunívoca entre clases de representaciones y clases de módulos de representación, por tanto con clases de grupos con operadores. La cuestión ulterior de la totalidad de las representaciones exige un estudio más preciso de la estructura de los anillos que han de representarse, y por tanto, en el caso especial, de los sistemas hipercomplejos. En el capítulo II se obtienen de nuevo los resultados de Wedderburn y se los lleva más lejos, sobre la base de la concepción teórico-grupal según la cual los ideales de un solo lado ocupan el primer plano. Resulta que el ``teorema de la cadena de múltiplos'' para ideales derechos, o la idéntica ``condición mínima'' (en cada conjunto de ideales derechos hay al menos un miembro minimal dentro del conjunto), basta como condición de finitud.\footnote{Los métodos de prueba de Wedderburn pueden transferirse si se supone el ``teorema de doble cadena''. Cf. E. Artin, \emph{Zur Theorie der hyperkomplexen Zahlen}, Hamb. Abh. 1927. Cf. también A. Suschkewitsch, \emph{Über die endlichen Gruppen ohne das Gesetz der eindeutigen Umkehrbarkeit}, Math. Annalen 99 (1928), pp. 30--50, donde se desarrollan teoremas estructurales paralelos para dominios (finitos) con una sola operación asociativa no invertible. La escisión del ``núcleo'' de Suschkewitsch en grupos derechos e izquierdos corresponde a la representación de un anillo bilateralmente simple con elemento identidad como suma de ideales derechos e izquierdos.} Se obtiene la identidad de los anillos sin radical que satisfacen la condición mínima con los anillos (derechamente) completamente reducibles con elemento identidad. En tales anillos todos los ideales derechos simples de un anillo bilateralmente indescomponible pertenecen a la misma clase y por ello poseen -- como grupos con operadores -- el mismo anillo de automorfismos, que por la simplicidad de los ideales se convierte en un anillo de división. Este anillo de división de automorfismos de la clase es el que media la representación matricial hallada por Wedderburn.

Esta representación matricial en el anillo de división de automorfismos resuelve también el problema de la representación en el caso completamente reducible, tan pronto como se exige la representación respecto del anillo de división de automorfismos o respecto de sus subcuerpos; en particular, respecto del dominio de coeficientes de un sistema hipercomplejo. No hay -- esto es consecuencia directa del teorema que reduce las clases de módulos a clases de ideales -- más representaciones irreducibles que la representación de Wedderburn y aquellas que se derivan de ella cuando los elementos del anillo de división de automorfismos se sustituyen de manera natural por matrices sobre un subcuerpo.

Así, en el caso completamente reducible hay tantas clases distintas de representaciones irreducibles como clases de ideales; su número coincide con el número de anillos distintos bilateralmente indescomponibles, y por tanto simples. Este número es finalmente de nuevo idéntico al número de componentes indescomponibles del centro, que se convierten en cuerpos conmutativos. En el caso de sistemas hipercomplejos con dominio de coeficientes algebraicamente cerrado, estos últimos cuerpos están completamente determinados por los distintos homomorfismos de anillos (representaciones irreducibles) del centro; salvo por un factor numérico, estos son precisamente los caracteres.

Los resultados obtenidos dan al mismo tiempo una visión de conjunto de las representaciones irreducibles de sistemas con radical, en la medida en que estas pueden reducirse a las representaciones del anillo de clases residuales módulo el radical, que pasa a ser un sistema sin radical.

Como se mostrará más adelante, para sistemas hipercomplejos sin radical la teoría de los ``cuerpos de descomposición'' se sigue también del anillo de división de automorfismos: es decir, la teoría de las extensiones conmutativas del dominio de coeficientes en las cuales las representaciones se descomponen en representaciones absolutamente irreducibles; equivalentemente, en las cuales se produce una descomposición en suma directa en ideales derechos absolutamente simples. Los cuerpos de descomposición son idénticos a los del anillo de división de automorfismos, y todos los cuerpos de descomposición de grado mínimo son isomorfos a los subcuerpos conmutativos máximos del anillo de división de automorfismos. La conexión con las investigaciones de Schur mencionadas arriba se establece por el hecho de que las transpuestas de las matrices que conmutan con la representación proporcionan precisamente una representación del anillo de división de automorfismos.\footnote{Cf. R. Brauer--E. Noether, \emph{Über minimale Zerfällungskörper irreduzibler Darstellungen}, Sitz.-Ber. Preuß. Akad. Wiss. 1927, p. 221.} Estos teoremas han de desarrollarse sistemáticamente en el marco de una teoría de Galois de los anillos de división no conmutativos.

Los conceptos subyacentes -- homomorfismo por operadores y anillo de automorfismos -- dan también la estructura de los anillos generales con radical que satisfacen solo la condición mínima; esto se llevará a cabo desde otro punto de vista.

\section*{Capítulo I. Fundamentos teórico-grupales}

\subsection*{§ 1. Grupos con operadores}

Sea $\mG$ un grupo (finito o infinito), con elementos $a,b,\ldots$. Por un dominio de operadores $\Omega$ para el grupo $\mG$ se entiende un conjunto de nuevos símbolos $H,\Theta,\ldots$, tal que a cada $a$ de $\mG$ y a cada $\Theta$ de $\Omega$ se asigna un $\Theta a$ unívocamente determinado en $\mG$, y tal que vale la ley distributiva:
\[
        \Theta(ab)=\Theta a\cdot\Theta b.
\]
Por consiguiente, cada operador define un homomorfismo del grupo en sí mismo, donde el grupo se aplica o bien sobre sí mismo o bien sobre un subgrupo. El elemento identidad pasa al elemento identidad, y los inversos pasan a inversos.

Un dominio de operadores se llama \emph{absoluto} si operadores distintos definen también homomorfismos distintos. De un dominio de operadores arbitrario se obtiene uno absoluto identificando todos aquellos elementos que generan el mismo homomorfismo. Un dominio de operadores absoluto es una imagen biunívoca de un subconjunto del conjunto de todos los homomorfismos del grupo en sí mismo.

Un \emph{subgrupo admisible} $\mH$ de $\mG$ -- respecto de un dominio de operadores fijo $\Omega$ -- es un subgrupo que admite el dominio de operadores, es decir, tal que $\Theta a$ está en $\mH$ para todo $a$ en $\mH$ y todo $\Theta$ en $\Omega$.

\emph{Ejemplos.} 1. Sean los operadores los automorfismos interiores: $\Theta a=c^{-1}ac$. Los subgrupos admisibles son los divisores normales.

2. Sean los operadores todos los automorfismos. Los subgrupos admisibles son los ``subgrupos característicos'', que van a sí mismos bajo todos los automorfismos.\footnote{Los conceptos explicados aquí proceden de W. Krull, \emph{Über verallgemeinerte endliche Abelsche Gruppen}, Math. Zeitschr. 23 (1925), pp. 161--196, y O. Schmidt, \emph{Über unendliche Gruppen mit endlicher Kette}, Math. Zeitschr. 29 (1928), pp. 34--41.}

3. Sea $\mG$ un anillo, es decir, un grupo abeliano respecto de la adición en el que además está definida una multiplicación, con las propiedades
\[
        r(a+b)=ra+rb,\qquad (a+b)r=ar+br,\qquad ab\cdot c=a\cdot bc.
\]
Cada elemento $r$ define simultáneamente dos operadores: los operadores $rx$ y $xr$. Los subgrupos admitidos son los ``ideales'' $\mathfrak a$, a saber: ideales izquierdos, que admiten las operaciones $rx$: $r\mathfrak a\subseteq\mathfrak a$; ideales derechos, que admiten las operaciones $xr$: $\mathfrak a r\subseteq\mathfrak a$; e ideales bilaterales, que admiten ambas operaciones. Todos los ideales se vuelven triviales $(=\{0\}\text{ o }=\mG)$ cuando el anillo $\mG$ es un anillo de división, es decir, cuando $\mG-\{0\}$ es un grupo respecto de la multiplicación.\footnote{Así, ni para anillos ni para anillos de división se supone la ley conmutativa de la multiplicación.}

4. \emph{Módulos respecto de un anillo $\mo$.} Sea $\mo$ un anillo y $\mM$ un grupo abeliano escrito aditivamente. Supóngase dada una multiplicación $r\cdot a$ de elementos de $\mo$ con elementos de $\mM$; el producto debe estar de nuevo en $\mM$. Se exige:
\[
\left.
\begin{aligned}
 r(a+b)&=ra+rb,\\
 rs\cdot a&=r\cdot sa,\\
 (r+s)a&=ra+sa,
\end{aligned}
\right\}\qquad r,s\in\mo,\quad a,b\in\mM.
\]
Entonces $\mM$ se llama un $\mo$-módulo; más precisamente, puesto que los multiplicadores de $\mo$ se escriben a la izquierda, un $\mo$-módulo izquierdo. El anillo $\mo$ es al mismo tiempo un dominio de operadores. Si es absoluto, es decir, si elementos distintos del anillo dan también operaciones distintas, se habla de un \emph{dominio absoluto de multiplicadores} para el módulo $\mM$.

Como antes, de un dominio de multiplicadores arbitrario se puede pasar al absoluto identificando los elementos que generan la misma operación. El dominio absoluto de multiplicadores es de nuevo un anillo. Los subgrupos admisibles de un módulo $\mM$ se llaman \emph{submódulos}. Si en particular $\mo=\mM$, se recuperan los ideales izquierdos.\footnote{De la misma manera se exige para módulos derechos:
\[
(a+b)r=ar+br,\qquad a\cdot rs=ar\cdot s,\qquad a(r+s)=ar+as.
\]}

5. \emph{Bimódulos.} Si $\mM$ es simultáneamente un $\mo$-módulo izquierdo y un $\mo'$-módulo derecho, y si además para $a$ en $\mo$, $\alpha$ en $\mM$ y $a'$ en $\mo'$ se tiene siempre
\[
        a\cdot\alpha a'=a\alpha\cdot a',
\]
entonces $\mM$ se llama un bimódulo.

6. \emph{El anillo de automorfismos de un grupo abeliano.}\footnote{Cf. A. Châtelet, \emph{Les Groupes Abéliens finis}, París, Gauthier-Villars, 1925, p. 99.} Para los endomorfismos de un grupo abeliano (escrito aditivamente) se puede definir una adición y una multiplicación mediante las fórmulas
\[
        (H+\Theta)a=Ha+\Theta a,\qquad (H\Theta)a=H(\Theta a).
\]
Las operaciones así definidas representan claramente de nuevo homomorfismos, y por tanto pertenecen otra vez al sistema. El sistema forma un anillo, y el grupo abeliano puede considerarse, según 4, como un módulo respecto de este anillo.

En lo que sigue, por ``subgrupos'' entenderemos siempre subgrupos admisibles. La intersección $\mA\cap\mB$ de dos subgrupos admisibles es de nuevo un subgrupo admisible. También lo es el producto $\mA\mB$, siempre que al menos uno de los dos subgrupos $\mA,\mB$ sea un divisor normal.

\subsection*{§ 2. Los teoremas de isomorfía}

Una aplicación de un grupo $\mG$ sobre un grupo $\overline{\mG}$ -- donde $\mG$ y $\overline{\mG}$ han de poseer el mismo dominio de operadores -- se llama un homomorfismo por operadores si, primero, es un homomorfismo en el sentido ordinario ($ab\mapsto \bar a\bar b$), y si, segundo, siempre que $a$ se aplica en $\bar a$, el elemento $\Theta a$ se aplica en $\Theta\bar a$.\footnote{Se puede extender el concepto de homomorfismo por operadores permitiendo que los grupos $\mG,\overline{\mG}$ tengan dominios de operadores separados; el homomorfismo aplica entonces no solo $\mG$ sobre $\overline{\mG}$, sino también los dominios de operadores entre sí, de tal manera que, si $a$ se aplica en $\bar a$ y $\Theta$ en $\overline\Theta$, entonces $\overline\Theta\bar a$ sea la imagen de $\Theta a$. En esta forma general, el concepto de homomorfismo por operadores incluye también el de homomorfismo de anillos; cf. §8.} Notación: $\mG\sim\overline{\mG}$. Si la correspondencia es biunívoca, se llama isomorfía por operadores. Notación: $\mG\simeq\overline{\mG}$. Como para los grupos ordinarios, se prueba el teorema de homomorfía:

\emph{Si $\overline{\mG}$ es una imagen operador-homomorfa de $\mG$, entonces $\overline{\mG}$ es operador-isomorfa a un grupo factor $\mG/\mN$, donde $\mN$ es un divisor normal admisible, formado por todos los elementos de $\mG$ que corresponden a la identidad de $\overline{\mG}$. Recíprocamente, todo divisor normal admisible $\mN$ define un grupo factor $\mG/\mN$ que admite el dominio de operadores de $\mG$ y es una imagen operador-homomorfa de $\mG$.}

Un grupo se llama \emph{simple} si no tiene divisores normales distintos de él mismo y del grupo identidad. Todo homomorfismo de un grupo simple es o bien una isomorfía, o bien asigna el elemento identidad a todos los elementos.

\emph{Primer teorema de isomorfía.} \emph{Sea $\overline{\mG}$ una imagen homomorfa de $\mG$, sea $\overline{\mA}$ un divisor normal de $\overline{\mG}$, y sea $\mA$ la totalidad de los elementos de $\mG$ a los que se asignan elementos de $\overline{\mA}$. Entonces $\mA$ es de nuevo un divisor normal, y}
\[
        \overline{\mG}/\overline{\mA}\simeq \mG/\mA. \tag{1}
\]
\emph{Demostración.} $\mG\sim\overline{\mG}$, $\overline{\mG}\sim\overline{\mG}/\overline{\mA}$, luego $\mG\sim\overline{\mG}/\overline{\mA}$. Por tanto $\overline{\mG}/\overline{\mA}$ es isomorfo a un grupo factor en $\mG$; el divisor normal correspondiente es la totalidad de los elementos que corresponden a la identidad en $\overline{\mG}/\overline{\mA}$; este es precisamente $\mA$.

\emph{Añadido.} Si, de acuerdo con el teorema de homomorfía, se pone $\overline{\mG}=\mG/\mN$, entonces $\mA$ contiene ciertamente a $\mN$. A partir de $\mA$ se puede recuperar $\overline{\mA}$: $\overline{\mA}=\mA/\mN$. La asignación $\mA\leftrightarrow\overline{\mA}$ es biunívoca. La fórmula (1) también puede escribirse
\[
        (\mG/\mN)/(\mA/\mN)\simeq \mG/\mA.
\]

\emph{Segundo teorema de isomorfía.} \emph{Sea $\mA$ un subgrupo y $\mB$ un divisor normal de $\mG$. Entonces $\mA\cap\mB$ es un divisor normal en $\mA$, y}
\[
        \mA\mB/\mB\simeq \mA/(\mA\cap\mB).
\]
\emph{Demostración.} Bajo el homomorfismo $\mG\sim\mG/\mB$, a los elementos $a$ de $\mA$ en particular se les asignan ciertas clases residuales $a\mB$, que juntas constituyen el grupo $\mA\mB/\mB$. Así $\mA\sim\mA\mB/\mB$, de donde la afirmación se sigue por el teorema de homomorfía.

Una aplicación operador-homomorfa de un grupo $\mG$ sobre sí mismo o sobre un subgrupo se llama endomorfismo de $\mG$. Si $\mG$ es en particular un grupo abeliano (con operadores), escrito aditivamente, y si la suma y el producto de homomorfismos se definen como arriba (§1, ejemplo 6), entonces los homomorfismos por operadores del grupo en sí mismo forman un anillo, el anillo de automorfismos.

\emph{El anillo de automorfismos de un grupo abeliano simple (con operadores) es un anillo de división.}

\emph{Demostración.} Todo homomorfismo aplica el grupo o bien sobre sí mismo o bien sobre el grupo cero (pues no hay otros subgrupos admisibles). Los homomorfismos que no envían todo a cero son, por lo observado arriba acerca de los grupos simples, isomorfías, y las aplicaciones isomorfas de un grupo sobre sí mismo forman un grupo. Así, después de omitir el operador cero, el anillo se convierte en un grupo; por tanto el anillo mismo es un anillo de división.

Si, en particular, $\mG$ es un $\mo$-módulo izquierdo, y si los homomorfismos por operadores se escriben como operadores derechos, entonces $\mG$ se convierte en un \emph{bimódulo respecto de $\mo$ a la izquierda y del anillo de automorfismos a la derecha}. Pues el hecho de que los nuevos operadores $\Gamma$ fueran elegidos como homomorfismos se expresa por las fórmulas
\[
        (a+b)\Gamma=a\Gamma+b\Gamma,
        \qquad ra\cdot\Gamma=r\cdot a\Gamma,
\]
que (junto con las restantes fórmulas ya interpretadas antes) caracterizan al bimódulo.

Recíprocamente, si se da un bimódulo, esas mismas fórmulas expresan que cada elemento del dominio de multiplicadores derechos induce un homomorfismo por operadores respecto de los operadores izquierdos, y recíprocamente.

\subsection*{§ 3. Series de composición}

Si en $\mG$ hay una sucesión finita de subgrupos admisibles
\[
        \mG>\mA_1>\mA_2>\cdots>\mA_r=\mE \tag{2}
\]
($\mE$ es el grupo que consiste solo en el elemento identidad $e$), tal que cada $\mA_i$ es un divisor normal en el miembro precedente, y es imposible insertar entre dos miembros sucesivos de la sucesión otro subgrupo con la misma propiedad, entonces la sucesión se llama una serie de composición.

Los grupos factor $\mG/\mA_1,\ \mA_1/\mA_2,\ldots,\ \mA_{r-1}/\mE$ se llaman factores de composición. Son grupos simples, es decir, no poseen divisores normales admisibles aparte de ellos mismos y del grupo identidad.

Si entre los operadores se incluyen, en particular, todos los automorfismos interiores, entonces la sucesión (2) consta enteramente de divisores normales y se la llama una serie principal. Si se incluyen todos los automorfismos, se llama una serie característica. En los ejemplos posteriores de §1 se trataría de series de composición de ideales en $\mo$, respectivamente series de composición de módulos o bimódulos.

\emph{Teorema de Jordan--Hölder.} \emph{Si para un grupo $\mG$ existen dos series de composición distintas,}
\[
        \mG>\mA_1>\mA_2>\cdots>\mA_r=\mE,
        \qquad
        \mG>\mB_1>\mB_2>\cdots>\mB_s=\mE,
\]
\emph{entonces tienen la misma longitud, $r=s$, y los factores de composición}
\[
        \mG/\mA_1,\ \mA_1/\mA_2,\ldots,\mA_{r-1}/\mE
\]
\emph{son, en algún orden, isomorfos a}
\[
        \mG/\mB_1,\ \mB_1/\mB_2,\ldots,\mB_{s-1}/\mE.
\]
Para la demostración véase, por ejemplo, E. Noether, \emph{Abstrakter Aufbau} usw., Math. Annalen 96 (1926), §10, p. 57. Allí se prueba al mismo tiempo que:

\emph{Si en $\mG$ hay una serie de composición, entonces puede trazarse una serie de composición a través de todo divisor normal $\mH$ de $\mG$.}

La existencia de una serie de composición es clara para grupos finitos y, más generalmente, para aquellos grupos en los que valen una condición maximal y una condición minimal:

La \emph{condición maximal} dice que en cada conjunto de subgrupos hay un subgrupo maximal, es decir, uno que ya no está contenido en otro subgrupo del conjunto. Equivalentemente, toda cadena ascendente de subgrupos $\mA_1\subset\mA_2\subset\mA_3\cdots$ se interrumpe después de un número finito de miembros.

La \emph{condición minimal} dice que en cada conjunto de subgrupos hay un subgrupo minimal, uno que no contiene ya a ningún otro subgrupo del conjunto; equivalentemente, toda cadena descendente $\mA_1\supset\mA_2\supset\mA_3\cdots$ se interrumpe después de un número finito de miembros.

\emph{Si solo se admiten divisores normales como subgrupos} (por ejemplo para grupos abelianos), \emph{entonces, recíprocamente, las condiciones maximal y minimal se siguen de la existencia de la serie de composición.}

\emph{Demostración.} Todo divisor normal posee una serie de composición, cuya longitud se llama longitud del divisor normal. Si $\mA\subset\mB$, entonces la longitud de $\mA$ es menor que la de $\mB$, ya que por $\mA$ puede trazarse una serie de composición para $\mB$. Por tanto, todo subgrupo de longitud mínima es al mismo tiempo minimal, y todo subgrupo de longitud máxima es al mismo tiempo maximal.

\subsection*{§ 4. Productos directos e intersecciones}

Un grupo $\mG$ se llama producto directo de dos factores, $\mG=\mA\times\mB$, si
\[
\begin{array}{ll}
1.& \mA,\mB\text{ son divisores normales en }\mG,\\
2.& \mA\mB=\mG,\\
3.& \mA\cap\mB=\mE.
\end{array}
\]
\emph{La definición es claramente equivalente a lo siguiente:} Todo elemento $g$ de $\mG$ tiene una representación única como $g=ab$, con $a$ en $\mA$ y $b$ en $\mB$, y los elementos de $\mA$ conmutan con los de $\mB$: $ab=ba$.

Un grupo $\mG$ se llama producto directo de $n$ factores, $\mG=\mA_1\times\cdots\times\mA_n$, si, poniendo $\mB_i=\mA_1\cdots\mA_{i-1}\mA_{i+1}\cdots\mA_n$, se tiene $\mG=\mA_i\times\mB_i$ directamente para cada $i$.

\emph{La definición es de nuevo equivalente a lo siguiente:} Todo $g$ de $\mG$ es representable de manera única como
\[
        g=a_1a_2\cdots a_n,\qquad a_i\in\mA_i,
\]
y los elementos de $\mA_i$ conmutan con los de $\mA_k$.

Los siguientes hechos se usan a menudo:
\begin{enumerate}
\item Si $\mA_1\times\cdots\times\mA_n=\mR$ y $\mR\times\mH=\mG$, entonces $\mA_1\times\cdots\times\mA_n\times\mH=\mG$.
\item Si $\mG=\mA_1\times\cdots\times\mA_n=\mC_1\times\cdots\times\mC_n$ y $\mC_i\subseteq\mA_i$, entonces $\mC_i=\mA_i$. (Esto se sigue de la representación de los elementos de $\mA_i$ mediante los $\mC_j$.)
\item Si $\mG=\mA\times\mB$ y $\mR\supseteq\mA$, entonces $\mR=\mA\times(\mR\cap\mB)$. (Esto se sigue de la representación de los elementos de $\mR$ en la forma $ab$, donde el segundo factor pertenece tanto a $\mR$ como a $\mB$.)
\item Si $\mG=\mA\times\mB$, entonces $\mA\sim\mG/\mB$ y $\mB\sim\mG/\mA$ (segundo teorema de isomorfía). De aquí se sigue además que:
\item Si $\mG=\mA\times\mB$, y si $\mA$ y $\mB$ poseen series de composición de longitudes $m$ y $n$, entonces $\mG$ posee una serie de composición de longitud $m+n$.
\item Si $\mA\sim\overline{\mA}$ y $\mB\sim\overline{\mB}$, entonces $\mA\times\mB\sim\overline{\mA}\times\overline{\mB}$.
\end{enumerate}
Un grupo se llama directamente indescomponible si no puede representarse como producto directo de factores $\ne\mE$.

Es claro que \emph{todo grupo que satisface la condición minimal es el producto directo de un número finito de grupos directamente indescomponibles}.\footnote{Como W. Krull en el caso conmutativo y O. Schmidt en general han mostrado (cf. nota 6), bajo la hipótesis de las condiciones maximal y minimal la representación es única salvo isomorfía. Puesto que se pueden agregar todos los automorfismos interiores al dominio de operadores sin cambiar el concepto de indescomponibilidad directa, basta exigir las condiciones de finitud para divisores normales, o la existencia de una serie principal.}

Si los grupos de que se trata se escriben aditivamente, entonces las nociones de producto y producto directo pasan a suma y suma directa. Notación: para una suma de grupos, $(\mA_1,\mA_2,\ldots)$; para una suma directa, $\mA_1+\mA_2+\cdots$.

El concepto de ``intersección directa'' no se usará más adelante, pero los siguientes teoremas acerca de la conexión entre intersección directa y producto directo muestran cómo la teoría de ideales de los próximos capítulos puede formularse también con intersecciones en lugar de sumas, restableciendo así la conexión con las teorías conmutativas conocidas.

Un grupo $\mD$ se llama intersección directa de dos grupos $\mR$ y $\mS$ respecto de $\mG$ si
\[
\begin{array}{ll}
1.& \mR,\mS\text{ son divisores normales en }\mG,\\
2.& \mR\cap\mS=\mD,\\
3.& \mR\mS=\mG.
\end{array}
\]
Análogamente, $\mD$ se llama intersección directa de $n$ grupos $\mS_1,\ldots,\mS_n$ si, poniendo $\mR_i=\mS_1\cap\cdots\cap\mS_{i-1}\cap\mS_{i+1}\cap\cdots\cap\mS_n$, la intersección $\mD=\mR_i\cap\mS_i$ es directa para cada $i$.

De la definición se sigue que $\mD$ debe ser un divisor normal en $\mG$. Si $\mD=\mS_1\cap\cdots\cap\mS_n$ es directa y, bajo el homomorfismo $\mG\sim\mG/\mD$, los grupos $\mG,\mS,\mD$ pasan a $\overline{\mG},\overline{\mS},\mE$, entonces $\mE=\overline{\mS}_1\cap\cdots\cap\overline{\mS}_n$ se vuelve directa. Recíprocamente, de cada representación de este tipo $\mE=\overline{\mS}_1\cap\cdots\cap\overline{\mS}_n$ se vuelve a una representación directa $\mD=\mS_1\cap\cdots\cap\mS_n$. Mediante esta asignación, los teoremas sobre intersecciones directas para $\mD$ arbitrario se reducen al caso especial $\mD=\mE$.

En el caso $\mD=\mE$ y dos factores, las condiciones de producto directo e intersección directa coinciden. Para $n$ factores existe la siguiente relación biunívoca entre producto directo e intersección:

\emph{I. Si $\mG=\mA_1\times\cdots\times\mA_n=\mA_i\times\mB_i$, entonces $\mE=\mB_1\cap\cdots\cap\mB_n=\mB_i\cap\mC_i$ directamente, y $\mC_i=\mA_i$.}

\emph{II. Si $\mE=\mS_1\cap\cdots\cap\mS_n=\mR_i\cap\mS_i$ directamente, entonces $\mG=\mR_1\times\cdots\times\mR_n=\mR_i\times\mT_i$, y $\mT_i=\mS_i$.}

\emph{Demostración de I.} Sea $\mB=\mB_1\cap\cdots\cap\mB_n=\mB_i\cap\mC_i$; entonces $\mC_i\supseteq\mA_i$. Mostramos que $\mC_i\subseteq\mA_i$. Si, por ejemplo, $c$ está en $\mC_1$, entonces $c$ está en $\mB_2,\ldots,\mB_n$, de modo que la representación por componentes respecto de los $\mA$ da
\[
        c=a_1a_2\cdots a_n=a_1ea_3\cdots a_n=a_1a_2e\cdots a_n=\cdots=a_1\cdots a_{n-1}e.
\]
Como el producto de todos los $\mA_i$ es directo, las representaciones coinciden; en toda posición salvo la primera aparece una vez un $e$, luego $c=a_1$, o $\mC_1\subseteq\mA_1$, es decir, $\mC_1=\mA_1$, y correspondientemente $\mC_i=\mA_i$. Se sigue que $\mB=\mB_i\cap\mC_i=\mB_i\cap\mA_i=\mE$; además $\mB_i\mC_i=\mB_i\mA_i=\mG$, de modo que la intersección es directa.

\emph{Demostración de II.} Sea $\mH=\mR_1\cdots\mR_n=\mR_i\mT_i$; entonces $\mT_i\subseteq\mS_i$. Mostramos que $\mS_i\subseteq\mT_i$. Sea, por ejemplo, $s$ en $\mS_1$. Mediante $\mG=\mS_i\times\mR_i$ se obtiene
\[
        s=s_1e=s_2r_2=\cdots=s_nr_n.
\]
Formando $t_1=er_2\cdots r_n=t_ir_i$, y teniendo en cuenta la conmutatividad elemento por elemento de $\mR_i$ y $\mS_i$, se obtiene
\[
        st_1^{-1}=s_it_i^{-1},
\]
que es por tanto elemento de todos los $\mS_i$, y de ahí igual a $e$. Así $\mS_1\subseteq\mT_1$, o bien $\mT_1=\mS_1$, y correspondientemente $\mT_i=\mS_i$. Se sigue que $\mH=\mR_i\mT_i=\mR_i\mS_i=\mG$ y $\mR_i\cap\mT_i=\mR_i\cap\mS_i=\mE$; por consiguiente, $\mG$ es el producto directo de los $\mR_i$.

\end{document}
